\documentclass[a4paper,11pt]{article}

\usepackage[margin=2.5cm]{geometry}
\usepackage[T1]{fontenc}
\usepackage[utf8]{inputenc}
\usepackage[english]{babel}
\usepackage{lmodern}
\usepackage{microtype}
\usepackage{amsmath}
\usepackage{booktabs}
\usepackage{tabularx}
\usepackage{float}
\usepackage{hyperref}
\usepackage{xcolor}

\hypersetup{
  colorlinks=true,
  linkcolor=blue!50!black,
  urlcolor=blue!50!black,
  pdftitle={IoT Challenge 3 -- Part 2 (Exercise)},
  pdfauthor={Tommaso Neri}
}

\setlength{\parindent}{0pt}
\setlength{\parskip}{0.55em}
\renewcommand{\arraystretch}{1.18}
\linespread{1.04}

\newcommand{\coursename}{Internet of Things Lab}
\newcommand{\challengename}{IoT Challenge \#3 -- Part 2 (Exercise)}
\newcommand{\authorname}{Tommaso Neri}
\newcommand{\personcode}{10817232}
\newcommand{\groupinfo}{Individual submission}
\newcommand{\answer}[1]{\textbf{\textcolor{red}{#1}}}

\title{\challengename}
\author{
  \authorname\\
  Person Code: \personcode\\
  \groupinfo
}
\date{\today}

\begin{document}

\maketitle
\thispagestyle{empty}

\section*{Scenario and Parameters}

The LoRaWAN network operates in Europe at carrier frequency $868$ MHz with
bandwidth $125$ kHz. It has one gateway and:
\[
  N = 40
\]
sensor nodes. Each node transmits packets with payload size 20 bytes according
to a Poisson process with intensity:
\[
  \lambda = 2\ \mathrm{packets/minute}
  = \frac{2}{60}\ \mathrm{packets/second}
  = \frac{1}{30}\ \mathrm{packets/second}.
\]

All nodes are assumed to use the same spreading factor. Following the simplified
ALOHA model used in the course material, the packet success rate is:
\[
  P_{\mathrm{success}} =
  \exp\left(-2N\lambda T_{\mathrm{air}}\right),
\]
where $T_{\mathrm{air}}$ is the packet airtime in seconds.

\section*{EQ1: Largest SF with at Least 75\% Packet Success Rate}

The target condition is:
\[
  \exp\left(-2N\lambda T_{\mathrm{air}}\right) \geq 0.75.
\]

Solving it for $T_{\mathrm{air}}$ gives:
\[
\begin{aligned}
T_{\mathrm{air}}
  &\leq \frac{-\ln(0.75)}{2N\lambda} \\
  &= \frac{-\ln(0.75)}{2 \cdot 40 \cdot \frac{1}{30}} \\
  &= 0.10788\ \mathrm{s} \\
  &= 107.88\ \mathrm{ms}.
\end{aligned}
\]

Therefore, the largest acceptable spreading factor is the largest SF whose
airtime is not greater than about $107.88$ ms.

\subsection*{Airtime Calculation}

The airtimes below are computed using the TTN LoRaWAN airtime calculator
(\url{https://www.thethingsnetwork.org/airtime-calculator}) with the parameters
specified by the assignment:

\begin{itemize}
    \item input bytes: 20;
    \item region: Europe 863--870 MHz;
    \item bandwidth: 125 kHz;
    \item spreading factor varied from SF7 to SF12.
\end{itemize}

The TTN calculator is a LoRaWAN airtime calculator: the 20 input bytes are the
application payload bytes, and the calculator includes the LoRaWAN overhead in
the time-on-air result. The carrier frequency does not affect airtime.

\begin{table}[H]
\centering
\begin{tabular}{ccc}
\toprule
SF & TTN time on air [ms] & ALOHA success rate \\
\midrule
7  & 71.9   & 82.54\% \\
8  & 133.6  & 70.02\% \\
9  & 246.8  & 51.78\% \\
10 & 452.6  & 29.91\% \\
11 & 987.1  & 7.19\% \\
12 & 1810.4 & 0.80\% \\
\bottomrule
\end{tabular}
\caption{TTN airtime and ALOHA success rate for 20 input bytes at BW 125 kHz.}
\end{table}

SF7 is below the maximum allowed airtime:
\[
  T_{\mathrm{air,SF7}} = 71.9\ \mathrm{ms} < 107.88\ \mathrm{ms},
\]
and its success rate is:
\[
  P_{\mathrm{success,SF7}}
  = \exp\left(-2 \cdot 40 \cdot \frac{1}{30} \cdot 0.0719\right)
  \approx 0.8254.
\]

SF8 is already above the airtime threshold when using the TTN LoRaWAN airtime
calculator:
\[
  T_{\mathrm{air,SF8}} = 133.6\ \mathrm{ms} > 107.88\ \mathrm{ms},
\]
and its success rate drops to about $70.02\%$.

\textbf{Answer EQ1:} \answer{Using the TTN LoRaWAN airtime calculator with 20 input bytes, the largest spreading factor that guarantees at least 75\% packet success rate is SF7.}

\section*{EQ2: Field Deployment with Non-Uniform Performance}

The measured packet success rate is significantly lower than expected, and the
performance is not uniform across nodes: some nodes work relatively well, while
others have very poor delivery rates.

The most appropriate action is:
\[
  \answer{\text{2. Move the nodes closer to the gateway.}}
\]

The non-uniformity across nodes suggests a coverage or link-budget problem.
Nodes that are farther from the gateway, or in worse propagation conditions,
may have received power close to or below the receiver sensitivity, causing poor
delivery even if the collision model predicted a better average success rate.

Decreasing the number of nodes would reduce the offered traffic and therefore
the collision probability, but it would mainly address a global congestion
problem. It does not directly explain why only some nodes perform very poorly.

Increasing the spreading factor can improve sensitivity and range, but it also
increases packet airtime. Longer airtime increases the probability of overlap
and collisions in the ALOHA model. Therefore, increasing SF is not the best
first action when the problem is non-uniform and likely related to node
placement or coverage.

\textbf{Answer EQ2:} \answer{Move the nodes closer to the gateway, because the observed non-uniform delivery rate points to coverage/link-budget issues rather than only network-wide congestion.}

\end{document}
